\scanpage{33}
trouver \(U\) \inserted{comme ci-dessus}, mais on peut trouver \(U\) de façon que l'on ait
\[
S_Y(U(a))=na
\]
alors, posant \(Z'(a)=Z(U(a))\), la relation
\[
Z'(S_Y(\nu))=nZ(\nu)
\]
\margin{\(\nu\) de degré 0 sur \(Y\).}
\deleted{Car on peut \unsure{écrire} \(nS(\nu)=U(a)\) où \(a=S_Y(\nu)\), \(S(\nu)\) de la forme \(U(a')\), car si \(a'=S_Y(\nu)\), on aura \(S_Y(U(a'))=na'=S_Y(n\nu)\) donc \(U(a')=n\nu\), et comme on aura \(a'=na\) pr \unsure{convenable} \(n\), il s'ensuivra \(U(a)\)}

Pour le voir, soit \(a=S_Y(\nu)\), on aura
\[
\deleted{\unsure{\operatorname{cl}Z'}}S_Y(U(a))=na=S_Y(n\nu)
\]
d'où \(\deleted{U(a)=n\nu,}\inserted{U(a)\unsure{=|\equiv}n\nu}\)
d'autre part on aura
\[
\begin{array}{ccc}
\deleted{\unsure{\operatorname{cl}}}Z'(S_Y(U(a)))&=&nZ(U(a))\\
\Vert&&\Vert\\
Z'(na)&&nZ'(a)
\end{array}
\]
\deleted{d'où, comme \(U(a)=n\nu\)}
\[
\deleted{nZ'(S_Y(\nu))=n^2Z(\nu)}
\]
Or tout élément de \(A_{0\,\mathrm{alb}}(Y)\) de degré 0 est de la forme \(n\unsure{\zeta}\) (car \(A_Y(k)\) est divisible!) donc appliquant ce qui précède à \(\unsure{\zeta}\) au lieu de \(\nu\), on trouve
\[
Z'(S_Y(\nu))=nZ(\nu).
\]
On prend comme isogénie la multiplication pr \(n\).
\mspage{32}
